\documentclass[11pt,a4paper]{article}
\usepackage[utf8]{inputenc}
\usepackage[T1]{fontenc}
\usepackage{amsmath,amssymb}
\usepackage[margin=2.4cm]{geometry}
\usepackage{hyperref}
\hypersetup{colorlinks=true,urlcolor=blue,linkcolor=blue}
\renewcommand{\abstractname}{Resumen}
\renewcommand{\contentsname}{Índice}
\setlength{\parindent}{0pt}
\setlength{\parskip}{0.6em}
\newcommand{\R}{\mathbb R}
\newcommand{\N}{\mathbb N}
\newcommand{\Bool}{\{\mathrm F,\mathrm V\}}
\newcommand{\len}[1]{\lvert #1\rvert}
\newcommand{\pre}{^{\mathrm{pre}}}
\newcommand{\wpof}[2]{\operatorname{wp}(#1,#2)}
\title{Trabajo práctico 1\\Fondo Monetario Común\\Especificación y weakest precondition}
\author{Juan Pablo Otero \and Fernando Toralez}
\date{Algoritmos y Estructuras de Datos, UBA\\Primer cuatrimestre de 2024}
\begin{document}
\maketitle
\begin{abstract}
Resolución grupal recuperada de Overleaf, revisada posteriormente sobre la
versión 3 de la consigna (8 de abril de 2024). Se conservan el problema y su
autoría; las correcciones de contratos, alcance y demostración se describen
en el README y en el historial. Este documento no es una corrección docente.
\end{abstract}
\tableofcontents
\clearpage
\section{Modelo y notación}
Las secuencias se indexan desde cero. $s[i]$ sólo se utiliza con
$0\leq i<\len{s}$. Las condiciones que habilitan índices se evalúan antes
de las expresiones que los emplean. Una cláusula sobre el estado previo
usa el superíndice $\mathrm{pre}$.

Sean $n$ individuos y $h$ pasos temporales. En apuestas y pagos, la fila
$i$ tiene $m_i$ eventos posibles. No se necesita que todos los $m_i$
coincidan. En eventos, todas las filas tienen la misma longitud $h$:
son instantes comparables del mismo proceso.

Escribimos $\operatorname{Config}(r,c,a,p,e)$ cuando:
\begin{itemize}
\item $n=\len{r}>0$ y $\len{c}=\len{a}=\len{p}=\len{e}=n$;
      $c$ es una secuencia de booleanos y $r_i>0$;
\item para todo $i<n$, $m_i=\len{a_i}=\len{p_i}>0$;
      $0\leq a_{ij}\leq1$, $\sum_{j=0}^{m_i-1}a_{ij}=1$ y $p_{ij}>0$;
\item para todo $i<n$, $\len{e_i}=h$ y
      $e_{it}\in\N,\ 0\leq e_{it}<m_i$ para todo $t<h$.
\end{itemize}
Se admite $h=0$. Los recursos iniciales son positivos, pero apostar cero
al evento ocurrido puede llevar un recurso posterior a cero. Por ello,
la redistribución y los estados intermedios admiten recursos no negativos.
La positividad estricta de ambas apuestas de la sección de correctitud
se conserva: es una precondición explícita de ese algoritmo particular.

\section{Especificaciones}
\subsection{Redistribución de los frutos}
Entrada: $r\in\operatorname{seq}(\R)$, $c\in\operatorname{seq}(\Bool)$.
Salida: $z\in\operatorname{seq}(\R)$.
\begin{itemize}
\item Requiere: $\len{r}=\len{c}=n$ y $r_i\geq0$ para todo $i<n$.
\item Asegura: $\len{z}=n$; si $n=0$, $z$ es vacía. Si $n>0$:
\end{itemize}
\[
 F(r,c)=\frac{\sum_{j=0}^{n-1}[c_j=\mathrm V]\,r_j}{n},
 \qquad
 z_i=F(r,c)+
 \begin{cases}0&c_i=\mathrm V,\\r_i&c_i=\mathrm F.\end{cases}
\]
El caso vacío no evalúa la división. Se reparte entre todos, no sólo
entre cooperadores. Se cumple
$\sum_i z_i=\sum_i r_i$: el fondo sustituye los aportes de quienes cooperan
y quienes no cooperan conservan sus recursos, además de recibir su parte.

\begin{samepage}
\subsection{Trayectorias de los frutos individuales}
Procedimiento sin resultado adicional; sólo modifica $H$, una secuencia
de trayectorias.
\begin{itemize}
\item Requiere: cada fila inicial de $H$ tiene un elemento $r_i=H\pre[i][0]$,
      y vale $\operatorname{Config}(r,c,a,p,e)$.
\item Modifica: $H$.
\item Asegura: $\len{H}=n$, cada fila tiene longitud $h+1$,
      $H_{i0}=r_i$, y para todo $0\leq t<h$:
\end{itemize}
\[
 u_{it}=H_{it}\,a_{i,e_{it}}\,p_{i,e_{it}},
 \qquad H_{i,t+1}=\operatorname{Redistribuir}(u_{\cdot t},c)_i.
\]
Los pagos son multiplicativos. El vector $u_{\cdot t}$ se calcula completo
antes de redistribuir: no se mezcla un estado nuevo de un individuo con
el estado anterior de otro.

Denotamos por $\operatorname{Hist}(H;r,c,a,p,e)$ la conjunción de
dimensiones, estado inicial y recurrencia anteriores. Para una configuración
válida existe una única $H$, por inducción sobre $t$. Escribimos
$\Phi_i(r,c,a,p,e)=H_{ih}$ para su recurso final.
\end{samepage}

\subsection{Trayectoria extraña escalera}
Entrada: $s\in\operatorname{seq}(\R)$. Salida booleana $b$.
Requiere: $\mathrm{True}$.
Definimos, sólo para índices de la secuencia:
\[
 \operatorname{MaxLocal}(s,i)\iff
 (i=0\ \lor\ s_i>s_{i-1})\ \land\
 (i=\len{s}-1\ \lor\ s_i>s_{i+1}).
\]
Cada disyunción protege el acceso a su vecino. Los extremos comparan con
su único vecino; en longitud uno no hay vecinos y la condición es vacua.
Asegura:
\[
 b\iff\sum_{i=0}^{\len{s}-1}[\operatorname{MaxLocal}(s,i)]=1.
\]
En particular, vacío da falso, un elemento da verdadero y dos elementos
dan verdadero exactamente cuando son distintos. Las mesetas no producen
máximos estrictos en sus posiciones empatadas.

\subsection{Un individuo decide si cooperar}
Entradas $i,r,a,p,e$; entrada/salida $c$.
\begin{itemize}
\item Requiere: $\operatorname{Config}(r,c\pre,a,p,e)$ y $0\leq i<n$.
\item Modifica: $c$.
\item Asegura: $\len{c}=n$ y $c_j=c\pre_j$ para todo $j\neq i$.
      Además, para cada $q\in\Bool$:
\end{itemize}
\[
 \Phi_i(r,c,a,p,e)\ \geq\
 \Phi_i(r,c\pre[i\leftarrow q],a,p,e).
\]
Se compara el recurso final, no el inmediato ni el total del grupo.
Ambas opciones utilizan los mismos recursos iniciales y eventos.
Cualquier opción empatada es válida; el verificador conserva la previa.

\begin{samepage}
\subsection{Un individuo actualiza su apuesta}
Entradas $i,r,c,p,e$; entrada/salida $a$.
\begin{itemize}
\item Requiere: $\operatorname{Config}(r,c,a\pre,p,e)$ y $0\leq i<n$.
\item Modifica: $a$.
\item Asegura: $\operatorname{Config}(r,c,a,p,e)$, y
      $a_j=a\pre_j$ para todo $j\neq i$.
      Para cada vector $v$ de longitud $m_i$ con $v_j\geq0$ y
      $\sum_jv_j=1$:
\end{itemize}
\[
 \Phi_i(r,c,a,p,e)\ \geq\
 \Phi_i(r,c,a\pre[i\leftarrow v],p,e).
\]
No se permite variar apuestas ajenas. Los cuantificadores contienen
todas las referencias a sus variables: ninguna trayectoria alternativa
queda fuera de alcance.

El dominio de $v$ es un simplex cerrado y acotado. Para un horizonte
finito, cada recurso final es un polinomio continuo en sus componentes;
por tanto existe un máximo. Si se exigiera $v_j>0$, el máximo podría
no existir: con un individuo, un paso y evento cero, el recurso final
es proporcional a $v_0$ y el supremo se alcanza en $v_0=1$.
La consigna general permite proporciones en $[0,1]$; no se traslada
a ella la precondición estricta del ejercicio de correctitud.
\end{samepage}

\section{Correctitud del trabajo puramente individual}
\subsection{Programa, precondición y notación uniforme}
La firma usa \texttt{recurso}, pero la inicialización del pseudocódigo de
la versión 3 escribe \texttt{recursos}. Se normaliza ese identificador a
\texttt{recurso}; sin esta corrección el programa tiene una variable libre.
No cambia el cálculo. En esta sección, $e=\texttt{eventos}$ es una única
secuencia booleana, no la matriz de eventos del modelo general.
Los subíndices $c,s$ de $a$ y $p$ representan los campos
\texttt{apuesta.c}, \texttt{apuesta.s}, \texttt{pago.c} y \texttt{pago.s}.
\begin{verbatim}
res := recurso;
i := 0;
while i < |eventos| do
    if eventos[i] then
        res := (res * apuesta.c) * pago.c
    else
        res := (res * apuesta.s) * pago.s
    endif;
    i := i + 1
endwhile
\end{verbatim}
Sea $a_c+a_s=1$, $a_c,a_s,p_c,p_s,\rho>0$, donde $\rho$ representa
\texttt{recurso}. Llamamos $P$ a esta precondición, junto con el tipado.
Los parámetros de entrada no cambian. Para $0\leq k\leq h=\len{e}$:
\[
 C_k=\sum_{t=0}^{k-1}[e_t=\mathrm V],\quad
 S_k=\sum_{t=0}^{k-1}[e_t=\mathrm F],\quad
 G_k=(a_cp_c)^{C_k}(a_sp_s)^{S_k}.
\]
Las sumas del prefijo vacío valen cero y $G_0=1$. La postcondición es
$Q:\ res=\rho G_h$.

\subsection{Invariante e inicialización}
Tomamos
\[
 I:\ P\land 0\leq i\leq h\land res=\rho G_i.
\]
Se conserva $P$ en el invariante, no se lo elimina antes de usar sus
hipótesis de definición. La sustitución para inicializar da:
\[
 \wpof{res:=\rho;\ i:=0}{I}
 \equiv P\land 0\leq0\leq h\land\rho=\rho G_0.
\]
Los parámetros tipados, $h\geq0$ y $G_0=1$ demuestran
$P\Rightarrow\wpof{res:=\rho;\ i:=0}{I}$.

\subsection{Preservación y weakest precondition}
Sea $B\equiv i<h$ y $C$ el condicional. Primero sustituimos el incremento:
\[
 \wpof{i:=i+1}{I}=
 P\land 0\leq i+1\leq h\land res=\rho G_{i+1}.
\]
Bajo $I\land B$, el índice $e_i$ está definido y las operaciones aritméticas
también lo están. La WP de $C$ exige la condición de la rama elegida:
\[
 \begin{aligned}
 &(e_i=\mathrm V\Rightarrow
       P\land 0\leq i+1\leq h\land res(a_cp_c)=\rho G_{i+1})\\
 {}\land{}&
 (e_i=\mathrm F\Rightarrow
       P\land 0\leq i+1\leq h\land res(a_sp_s)=\rho G_{i+1}).
 \end{aligned}
\]
Si $e_i=\mathrm V$, $C_{i+1}=C_i+1$, $S_{i+1}=S_i$ y
$G_{i+1}=G_i(a_cp_c)$. Sustituyendo $res=\rho G_i$ se obtiene
$\,res(a_cp_c)=\rho G_{i+1}$. La rama falsa es análoga con
$S_{i+1}=S_i+1$ y $G_{i+1}=G_i(a_sp_s)$.
La guarda y la integralidad de $i$ garantizan $0\leq i+1\leq h$.
Por tanto:
\[
 I\land B\Rightarrow\wpof{C;\ i:=i+1}{I}.
\]
No se divide ni cancela $\rho$: desaparecía incorrectamente en varios
pasos de la demostración histórica. El factor inicial se mantiene en
ambas ramas y en la postcondición.

\subsection{Salida y terminación}
De $I\land\neg B$ se deduce $i=h$ y, por sustitución,
$res=\rho G_h$, que es $Q$. Esto acredita corrección parcial.

La variante entera $V=h-i$ es no negativa bajo $I$ y estrictamente positiva
bajo $I\land B$. El cuerpo no altera $h$ y el incremento produce
$V'=h-(i+1)=V-1<V$. No puede existir una cadena infinita decreciente de
enteros no negativos. El ciclo termina, por lo que la corrección es total.
El caso $h=0$ inicializa $res=\rho$ y sale sin consultar ningún evento.

\section{Comprobaciones y fuentes}
El verificador auxiliar usa racionales exactos para contrastar conservación
de recursos, trayectorias, máximos locales, elección entre cooperar o no
y la fórmula individual, incluyendo prefijos del invariante.
Una malla finita de apuestas sólo verifica ejemplos; no demuestra un
óptimo sobre todos los reales. Ese requisito está expresado por el
cuantificador y la existencia se justifica por continuidad y compacidad.

Consigna localizada en una copia pública del documento de la cátedra:
\href{https://github.com/SashaBerkowsky/AyED-II/blob/42722b4636a1e8e3550f9c6500d3e3ed61e67792/Trabajos\%20Pr\%C3\%A1cticos/TP1/enunciado.pdf}{Fondo Monetario Común, versión 3, 1C2024}.
No se incorporaron soluciones de ese repositorio.
La versión pública usa paquetes estándar y macros propias; el proyecto
original con auxiliares y logos permanece íntegro en el respaldo académico.
\end{document}
